\documentclass[10pt,a4paper]{amsart}
\usepackage[margin=19mm]{geometry}
\usepackage{amsmath,amssymb,mathtools}
\usepackage[scr=rsfs,cal=euler]{mathalfa}
\usepackage{xfrac}
\usepackage[unicode,hidelinks,bookmarks=false]{hyperref}
\newcommand{\ie}{i.e.\ }
\newcommand{\cf}{cf.\ }
\newcommand{\resp}{respectively\ }
\newcommand{\Subsection}{\subsection}
\providecommand{\nobreakdash}{\nobreak\hbox{-}}
\setlength{\emergencystretch}{2em}
\multlinegap0pt
\usepackage{latexsym}
\newtheorem{thrm}{Theorem}[section]
\newtheorem{lemma}[thrm]{Lemma}
\newtheorem{prop}[thrm]{Proposition}
\newtheorem{defn}[thrm]{Definition}
\newtheorem{cor}[thrm]{Corollary}
\newtheorem{remark}[thrm]{Remark}
\newtheorem{exam}{Example}
\newtheorem{ques}{Question}
\newtheorem{assumption}[thrm]{Assumption}
\numberwithin{equation}{section}
\usepackage{color}
\usepackage{mathtools}
\mathtoolsset{showonlyrefs}
\usepackage{mathrsfs}
\def\P{\mathbb{P} }
\def\Q{\mathbb{Q} }
\def\R{\mathbb{R} }
\def\N{\mathbb{N} }
\def\V{\mathbb{V} }
\def\D{\mathbb{D} }
\def\cD{\mathcal{D} }
\def\cA{\mathcal{A} }
\def\B{\mathcal{B} }
\def\bP{{\bf P} }
\def\cL{{\cal L} }
\def\C{{\cal C} }
\def\doublespace{\baselineskip=16pt}
\def\singlespace{\baselineskip=12pt}
\begin{document}
\title[On the maximal displacement of subcritical branching random ]{On the maximal displacement of subcritical branching random walks with stretched exponential tail}
\author{f̱ Haojie Hou; Haojie Hou}
\date{}
\maketitle
\noindent\emph{Source and licence.} On the maximal displacement of subcritical branching random walks with stretched exponential tail. Version arXiv:2606.28631v1 (CC BY 4.0). This material is distributed under the Creative Commons Attribution 4.0 International licence, http://creativecommons.org/licenses/by/4.0/, as recorded in the arXiv OAI licence metadata for this exact version and shown on the abstract page https://arxiv.org/abs/2606.28631v1. Full rights notice: licenses/arxiv-2606.28631-CC-BY-4.0.txt.\par\smallskip
\noindent\textit{Editorial scope.} Counted unit: the original \texttt{prop} environment \texttt{prop2} at source lines 614--627 together with its complete proof at lines 734--734; the delivered document keeps source lines 219--735 so that every referenced label and every citation resolves inside it. Native editable source is the paper's own concatenated TeX; the AMS wrapper is editorial formatting only. No publisher class, upstream script or unrelated artwork is redistributed.
 	 \begin{align}\label{Def-of-stopping-time}
 	 	S_n := \sum_{i=1}^n X_i \quad \text{and} \quad T_x^+ := \inf\left\{n \in \mathbb{N} : S_n > x\right\}.
 	 \end{align}
 	 The main goal of this section is to prove the following proposition.
 
	\begin{prop}\label{prop1}
		\begin{itemize}
			\item[(i)] For any $x > 1$ and $n \in \mathbb{N}$,
			\[
			m^{n/2} \, \mathbb{P}\left(T_x^+ = n \right) \leq m^{n/2} \, \mathbb{P}\left(S_n > x \right) \lesssim \mathbb{P}(X > x).
			\]
			\item[(ii)] For each $n \in \mathbb{N}$,
			\[
			\lim_{x \to +\infty} \frac{\mathbb{P}(T_x^+ = n)}{\mathbb{P}(X > x)} = 1.
			\]
		\end{itemize}
	\end{prop}
	
	We first prove Proposition \ref{prop1} for the case $b = 0$. Before that, we gather some useful results in the following lemma.
	
	\begin{lemma}\label{lem-useful-fact}
		Let $\ell$ be a slowly varying function.
		
		\begin{itemize}
			\item[(i)] For any $0<z_1\leq z_2$, 
			$
			\lim_{x \to +\infty} \sup_{ y\in [z_1, z_2]} \big|  \ell(xy)/\ell(x) - 1\big| =0.
			$
			
			\item[(ii)] For any $A > 1$ and $\delta > 0$, there exists $K(A, \delta)$ such that
			\[
			\frac{\ell(y)}{\ell(x)} \leq A \max\left\{(y/x)^\delta, (y/x)^{-\delta}\right\} \quad \text{for all } x, y \geq K(A, \delta).
			\]
			Consequently, for any $\delta > 0$, we have
			\[
			x^{-\delta} \lesssim_\delta \ell(x) \lesssim_\delta x^{\delta}.
			\]
		\end{itemize}
	\end{lemma}
	\textbf{Proof: }  For (i), see \cite[Theorem 1.5.2, p.22]{BGT1987}. (ii) is known as Potter's theorem and can be found in \cite[Theorem 1.5.6, p.25]{BGT1987}
	
	
	\hfill$\Box$
	

\noindent
\textbf{Proof of Proposition \ref{prop1} for $b=0$:}  (i) The first inequality holds trivally, so we only prove the second inequality. 
If $x\leq n$, then according to the inequality $\mathbb{P}(X>n)= \ell(n)n^a\gtrsim m^{n/2}$, we get 
\begin{align}\label{e1}
	m^{n/2} \mathbb{P}\left(S_n>x\right)\leq m^{n/2}\lesssim   \mathbb{P}(X>n)\leq   \mathbb{P}(X>x). 
\end{align}
If $x>n$, then combining the union bound and Lemma \ref{lem-useful-fact} (ii) (with $A=2, \delta=1$),  we have $\ell(x/n)/\ell(x)\lesssim n$ and 
\begin{align} \label{e2}
	& m^{n/2}  \mathbb{P}\left(S_n> x \right)  \leq nm^{n/2}  \mathbb{P}(X>x/n) = nm^{n/2} (x/n)^a \ell(x/n) \nonumber\\
	&= m^{n/2} n^{1-a} \frac{\ell(x/n)}{\ell(x)} \mathbb{P}(X>x) \lesssim m^{n/2} n^{2-a} \mathbb{P}(X>x) \nonumber\\
	&\lesssim \mathbb{P}(X>x) .
\end{align}
Combining \eqref{e1} and \eqref{e2} implies (i).

 (ii) Since $\{T_x^+=1\}= \{X_1>x\}$, the limit holds trivally. We consider the case $n\geq 2$ here. 
 For each $\varepsilon\in (1/2,1)$, define 
 \begin{align}\label{def-of-E-n}
 E_n:= \left\{ \mbox{there exists at most one }i \leq n \mbox{ such that } X_i>x^\varepsilon \right\}.
 \end{align}
 Fixing any $\Gamma>0$ with $\Gamma(2\varepsilon+1)<-(2\varepsilon-1)a$. Since $x^{-\Gamma}\lesssim_\Gamma \ell(x)\lesssim_\Gamma x^\Gamma$ by Lemma \ref{lem-useful-fact} (ii),  
 \begin{align}\label{e4'}
 & \mathbb{P}\left(E_n^c\right) =\mathbb{P}\left(\exists i\neq j \mbox{ such that } X_i, X_j>x^\varepsilon \right)\leq n^2 \mathbb{P}(X>x^\varepsilon)^2 \nonumber\\
 &= n^2\frac{\left(\ell(x^\varepsilon)\right)^2}{\ell(x)} x^{(2\varepsilon-1)a}  \mathbb{P}(X>x) \lesssim_\Gamma n^2\frac{x^{2\varepsilon \Gamma}}{x^{-\Gamma}} x^{(2\varepsilon-1)a}  \mathbb{P}(X>x) \nonumber\\
 & = n^{2} x^{(2\varepsilon+1)\Gamma +(2\varepsilon-1)a} \mathbb{P}(X>x). 
 \end{align}
 Thus,    $\lim_{x\to +\infty} \frac{\mathbb{P}(E_n^c)}{\mathbb{P}(X>x)} =0$. 
 Since for $x$ large enough such that $x/n>x^{\varepsilon}$,  we have
 \[
 \{T_x^+=n\}\subset \{S_n>x\}\subset \{\exists\ 1\leq j\leq n\ s.t.\ X_j>x/n \}\subset \{\exists\ 1\leq j\leq n\ s.t.\ X_j>x^{\varepsilon} \}.
 \]
 Therefore, on $\{T_x^+=n\}\cap E_n$, there exists a unique $j$ such that $X_j>x^\varepsilon$ and that $\max_{\ell\neq j} X_\ell \leq x^{\varepsilon}$, which together with $S_n >x$ implies that $X_j> x-nx^\varepsilon>x^\varepsilon$ for large $x$. In conclusion, we have  
\begin{align}\label{e4}
	& \lim_{x\to +\infty} \frac{\mathbb{P}(T_x^+ =n)}{\mathbb{P}(X>x)}   =\lim_{x\to +\infty} \frac{1}{\mathbb{P}(X>x)}\mathbb{P}\left(T_x^+=n, \exists !\ 1\leq j\leq n\ \mbox{such that } X_j>x^\varepsilon \right) \nonumber\\
	&= \lim_{x\to +\infty} \frac{1}{\mathbb{P}(X>x)} \mathbb{P}\left(T_x^+=n, \exists \ 1\leq j\leq n\ \mbox{such that } X_j>x- nx^\varepsilon, \max_{\ell\neq j} X_\ell \leq x^\varepsilon \right). 
\end{align}
Now we would like to replace $\max_{\ell\neq j} X_\ell \leq x^\varepsilon $ by $\max_{\ell\neq j} |X_\ell| \leq x^\varepsilon $. Noticing that 
\begin{align}\label{e11-3}
	 & \lim_{x\to +\infty} \frac{1}{\mathbb{P}(X>x)}\mathbb{P}\left(\exists \ 1\leq j\leq n\ \mbox{such that } X_j>x- nx^\varepsilon, \max_{\ell\neq j} |X_\ell| > x^\varepsilon \right) \nonumber\\
	 &\leq \lim_{x\to +\infty} \frac{n}{\mathbb{P}(X>x)} \mathbb{P}\left(X_1>x- nx^\varepsilon, \max_{\ell\neq 1} |X_\ell| > x^\varepsilon \right)\nonumber\\
	 &=  \lim_{x\to +\infty} \frac{n}{\mathbb{P}(X>x)}\mathbb{P}\left(X>x- nx^\varepsilon\right)\mathbb{P}\left(  |X_\ell| > x^\varepsilon \right)^{n-1} = n \lim_{x\to +\infty} \mathbb{P}\left(  |X_\ell| > x^\varepsilon \right)^{n-1}  =0,
\end{align}
where the second equality follows from  Lemma \ref{lem-useful-fact} (i). Plugging this back to \eqref{e4} yields that 
\begin{align}\label{e5}
	& \lim_{x\to +\infty} \frac{\mathbb{P}(T_x^+ =n)}{\mathbb{P}(X>x)} \nonumber\\
	&  = \lim_{x\to +\infty} \frac{1}{\mathbb{P}(X>x)} \mathbb{P}\left(T_x^+=n, \exists \ 1\leq j\leq n\ \mbox{such that } X_j>x- nx^\varepsilon, \max_{\ell\neq j} |X_\ell | \leq x^\varepsilon \right). 
\end{align}
If the unique $j$ with $X_j>x-nx^\varepsilon$ and $\max_{\ell\neq j} |X_\ell|\leq x^\varepsilon$ satisfies $j\leq n-1$, then on $T_x^+=n$, it must hold that $X_j = S_{n-1}- \sum_{\ell\neq j ,\ell\leq n-1} X_\ell \leq x+ (n-2)x^\varepsilon\leq x+nx^\varepsilon $.  Therefore, 
\begin{align}\label{e5'}
	& \limsup_{x\to +\infty} \frac{1}{\mathbb{P}(X>x)}\mathbb{P}\left(T_x^+=n, \exists \ 1\leq j\leq n-1\ \mbox{such that } X_j>x- nx^\varepsilon, \max_{\ell\neq j} |X_\ell | \leq x^\varepsilon \right) \nonumber\\
	& \leq \limsup_{x\to +\infty} \frac{1}{\mathbb{P}(X>x)} \mathbb{P}\left( \exists 1\leq j\leq n-1\ \mbox{such that }  |X_j-x|\leq nx^\varepsilon\right) \nonumber\\
	&\leq n  \limsup_{x\to +\infty} \frac{1}{\mathbb{P}(X>x)}\mathbb{P}\left(   |X-x|\leq nx^\varepsilon\right) \nonumber\\
	&= n  \lim_{x\to +\infty} \frac{\ell(x-nx^\varepsilon)(x-nx^\varepsilon)^a-\ell(x+nx^\varepsilon)(x+nx^\varepsilon)^a}{\ell(x)x^a}=0,
\end{align}
where in the last equality we used  Lemma \ref{lem-useful-fact} (i). Now plugging this back to \eqref{e5} and noticing that 
$\{X_n>x- nx^\varepsilon, \max_{\ell \leq n-1} |X_\ell | \leq x^\varepsilon \}\subset  \{ \max_{1\leq i\leq n-1} S_i\leq x\}$ and $\{X_n>x+ nx^\varepsilon, \max_{1\leq \ell \leq n-1} |X_\ell | \leq x^\varepsilon \}\subset  \{ S_n > x\}$  for large $x$, by Lemma \ref{lem-useful-fact} (i), we conclude that 
\begin{align}
	&1= \lim_{x\to +\infty} \frac{1}{\mathbb{P}(X>x)} \mathbb{P}\left(X>x+ nx^\varepsilon\right) \mathbb{P}\left( |X| \leq x^\varepsilon \right)^{n-1}  \nonumber\\
	&= \lim_{x\to +\infty} \frac{1}{\mathbb{P}(X>x)} \mathbb{P}\left(X_n>x+ nx^\varepsilon, \max_{1\leq \ell \leq n-1} |X_\ell | \leq x^\varepsilon \right) \nonumber\\
	&\leq  \lim_{x\to +\infty} \frac{\mathbb{P}(T_x^+ =n)}{\mathbb{P}(X>x)}   = \lim_{x\to +\infty} \frac{1}{\mathbb{P}(X>x)}\mathbb{P}\left(S_n >x, X_n>x- nx^\varepsilon, \max_{1\leq \ell \leq n-1} |X_\ell | \leq x^\varepsilon \right) \nonumber\\
	& \leq  \lim_{x\to +\infty} \frac{\mathbb{P}\left( X_n>x- nx^\varepsilon \right)}{\mathbb{P}(X>x)} =1,
\end{align}
which implies the desired result for $b=0$.

\hfill$\Box$



Now we prove Proposition \ref{prop1} for $b \in (0,1)$. The proof idea is quite similar to the case $b = 0$, but requires more delicate estimates. Observe that if $g(x)$ is a positive function such that $\lim_{x \to +\infty} \frac{g(x)}{x^{1-b}} = 0$, then by Taylor's expansion, uniformly for all $|y| \leq g(x)$, we have $\lim_{x \to +\infty} ((x+y)^b - x^b) = 0$. Therefore, together with Lemma \ref{lem-useful-fact}(i), we obtain that uniformly for all $|y| \leq g(x)$,
\begin{align}\label{proper-of-tail}
	\lim_{x \to +\infty} \frac{\mathbb{P}(X > x + y)}{\mathbb{P}(X > x)} = \lim_{x \to +\infty} \frac{\ell(x + y) (x + y)^a}{\ell(x) x^a} e^{-\lambda ((x+y)^b - x^b)} = 1.
\end{align}
Set $\mu := \mathbb{E}(X)$, $\sigma := \sqrt{\operatorname{Var}(X)}$, and $\widehat{X} := \frac{X - \mu}{\sigma}$. Then, under {\bf(H1)}, 
\[
\mathbb{P}(\widehat{X} > x) = \mathbb{P}(X > \sigma x + \mu) = \ell(\sigma x + \mu) (\sigma x + \mu)^a e^{-\lambda (\sigma x + \mu)^b} =: \widehat{\ell}(x) x^a e^{-\lambda \sigma^b x^b}
\]
for some slowly varying function $\widehat{\ell}$ and for all $x > (1 - \mu)/\sigma$. Since we assumed $\mathbb{E}(X^2) < \infty$ in {\bf(H1)}, by \cite[Theorem 8.3]{DDS2008} (with $R(x) = \lambda \sigma^b x^b$), uniformly for all $n < \frac{1}{2} (\log z)^3$,
\[
\lim_{z \to +\infty} \sup_{n<\frac{1}{2}(\log z)^3} \left| \frac{\mathbb{P}(S_n > n\mu + \sigma z)}{n \mathbb{P}(X > \sigma z + \mu)} -1\right|=0.
\]
Replacing $n\mu + \sigma z$ by $x$ and noting that $n < (\log x)^3$ whenever $n < \frac{1}{2} (\log z)^3$ for sufficiently large $x$, it follows from \eqref{proper-of-tail} that
\begin{align}\label{e6}
	\lim_{x \to +\infty} \sup_{n < (\log x)^3} \left| \frac{\mathbb{P}(S_n > x)}{n \mathbb{P}(X > x)} - 1 \right| = 0.
\end{align}
 
\begin{lemma}\label{lem2}
	For $x>1$, define $\varepsilon_0:= \frac{1}{2}\min\{ b, 1-b\}$ and
	\[
	\theta(x):= \lambda x^{b-1}-\frac{(\log x)^2}{x}.
	\]
	 Then there exists a constant $C_*>0$ such that for large $x$,
	 \[
	 \mathbb{E}\left(e^{ \theta(x) \min\{ X-\mu, x-x^{\varepsilon_0} \} }\right)  \leq \exp\left\{C_* \theta^2(x)\right\}.
	 \]
\end{lemma}
\textbf{Proof: }  Set $x_0:=x-x^{\varepsilon_0}$ for simplicity. Then 
\begin{align}
	 & \mathbb{E}\left( e^{ \theta(x) \min\{ X-\mu, x_0 \} }\right) -1 = e^{\theta(x)x_0}\mathbb{P}(X\geq x_0+\mu)+\mathbb{E}\left(e^{\theta(x)(X-\mu)}1_{\{X-\mu <x_0\}}\right)-1\nonumber\\
	 & = e^{\theta(x)x_0}\mathbb{P}(X\geq x_0+\mu)+ \mathbb{E}\left(e^{\theta(x)(X-\mu)}1_{\{(\theta(x))^{-1}< X-\mu <x_0\}}\right)\nonumber\\
	 &\quad +\left(\mathbb{E}\left(e^{\theta(x)(X-\mu)}1_{\{X-\mu \leq \theta(x)^{-1}\}}\right)-1\right)\nonumber\\
	 &=:  I_1+I_2+I_3.
\end{align}

\noindent
{\bf Estimate for $I_1$.} From $\varepsilon_0< 1-b$ we have  $\lim_{x\to+\infty} (x_0^b - x^b )=  0$. Therefore, 
 combining $\ell(x)\lesssim x$ and $\lim_{x\to+\infty} \theta(x)(x-x_0)=  0$, it holds that 
\begin{align}\label{bound-for-I1}
	&I_1 = e^{\theta(x)x_0}\ell(x_0+\mu) |x_0+\mu|^a e^{-\lambda (x_0+\mu)^b}\lesssim e^{\theta(x)x} x^{|a|+1}e^{-\lambda x^b}\nonumber\\
	& = e^{-(\log x)^2} x^{|a|+1}\lesssim \theta^2(x),
\end{align}
where in the last inequality we used the fact that $\theta(x)\gtrsim x^{b-1}$.

\noindent
{\bf Estimate for $I_2$.} Combining $\ell(x)\lesssim x$ and $\lceil x_0 \rceil \leq \lfloor x\rfloor$, we get
\begin{align}
	& I_2 \leq \sum_{k= \lfloor (\theta(x))^{-1}\rfloor }^{\lceil x_0\rceil} e^{\theta(x)(k+1)} \mathbb{P}\left( X-\mu\in (k,k+1]\right)\leq \sum_{k= \lfloor (\theta(x))^{-1}\rfloor}^{\lceil x_0\rceil} e^{\theta(x)(k+1)} \mathbb{P}\left( X-\mu >k\right)\nonumber\\
	&\lesssim \sum_{k= \lfloor (\theta(x))^{-1}\rfloor}^{ \lfloor x\rfloor} e^{\theta(x)k} k^{a+1} e^{-\lambda k^b} \lesssim \int_{(\theta(x))^{-1}/2}^{x} e^{\theta(x) y} y^{a+1}e^{-\lambda y^b} \mathrm{d}y. 
\end{align}
Noticing that $\theta (x)y -\lambda y^b\leq  \lambda x^{b-1}y -\lambda y^b \leq -\lambda y^b(1-2^{b-1})$ for $y<x/2$ and 
\begin{align}
\theta (x)y -\lambda y^b & =  \lambda x^{b-1}y^b(y^{1-b}-x^{1-b})  - (\log x)^2 y/x  \leq - (\log x)^2 /2
\end{align}
for $x/2\leq y<x$,   we conclude that 
\begin{align}\label{bound-for-I2}
	& I_2 \lesssim \int_{(\theta(x))^{-1}/2}^{x/2}  y^{a+1}e^{-\lambda y^b(1-2^{b-1})} \mathrm{d}y  +\int_{x/2}^{x}  y^{a+1}e^{-(\log x)^2/2 } \mathrm{d}y\nonumber\\
	&\lesssim \theta^2(x)\int_{(\theta(x))^{-1}/2}^{\infty}  y^{a+3}e^{-\lambda y^b(1-2^{b-1})} \mathrm{d}y  + x^{a+2}e^{-(\log x)^2/2 } \lesssim \theta^2(x).
\end{align}

\noindent
{\bf Estimate for $I_3$.} Combining inequalities 
\[
\mathbb{E}\left(\theta(x)(X-\mu)1_{\{X-\mu\leq \theta(x)^{-1}\}}\right)=-\mathbb{E}\left(\theta(x)(X-\mu)1_{\{X-\mu> \theta(x)^{-1}\}}\right) <0
\]
and  $|e^{x}-1-x|\lesssim x^2$ for $x\leq 1$, we get 
\begin{align}\label{bound-for-I3}
	& I_3 = \mathbb{E}\left(\left(e^{\theta(x)(X-\mu)} -1-\theta(x)(X-\mu)\right)1_{\{X-\mu \leq \theta(x)^{-1}\}}\right)-\mathbb{P}(\theta(x)(X-\mu)>1)\nonumber\\
	&\qquad + \mathbb{E}\left(\theta(x)(X-\mu)1_{\{X-\mu\leq \theta(x)^{-1}\}}\right)\nonumber\\
	&\leq \mathbb{E}\left(\left|e^{\theta(x)(X-\mu)} -1-\theta(x)(X-\mu)\right|1_{\{X-\mu \leq \theta(x)^{-1}\}}\right)\nonumber\\
	&\lesssim \theta^2(x) .
\end{align}
Thus, combining \eqref{bound-for-I1}, \eqref{bound-for-I2} and \eqref{bound-for-I3} implies the desired result. 


\hfill$\Box$

\noindent
\textbf{Proof of Proposition \ref{prop1} for $b\in (0,1)$:}
Let $\varepsilon_0$ be given as in Lemma \ref{lem2}. 

(i) We divide the proof into three cases.

\noindent
{\bf Case 1: $n<(\log x)^3$.}
 By \eqref{e6}, we have
  \begin{align}\label{e7}
		m^{n/2}\mathbb{P}(T_x^+=n)\leq m^{n/2} \mathbb{P}(S_n>x) \lesssim n m^{n/2} \mathbb{P}(X>x)\lesssim   \mathbb{P}(X>x).
\end{align}

\noindent
{\bf Case 2: $n\geq x^{(1+b)/2}$.}
Noticing that in this case, $n\log (1/m) /2 >2\lambda x^b$ when $x$ is large enough. Therefore, it follows from $\mathbb{P}(X>x)\gtrsim e^{-2\lambda x^b}$  that 
\begin{align}\label{e8}
	\mathbb{P}(X>x)\gtrsim e^{-2\lambda x^b} \geq m^{n/2} \quad  \Longrightarrow \quad m^{n/2}\mathbb{P}(T_x^+=n)\leq m^{n/2} \lesssim  	\mathbb{P}(X>x).
\end{align}

\noindent
{\bf Case 3: $(\log x)^3\leq n <x^{(1+b)/2}$.}
Let $x$ be large enough such that $x^{\varepsilon_0}>\mu$,  then by Markov inequality, 
\begin{align}\label{e33-2}
	&\mathbb{P}\left(S_n >x, \max_{1\leq i\leq n} X_i \leq x-2x^{\varepsilon_0} \right)  \leq \mathbb{P}\left(S_n-n\mu >x-n\mu, \max_{1\leq i\leq n} (X_i-\mu) \leq x-x^{\varepsilon_0} \right) \nonumber\\
	& \leq  \mathbb{P}\left(\sum_{i=1}^n \min\{ X_i-\mu, x-x^{\varepsilon_0}\}>x-n\mu \right)\nonumber\\
	& \leq e^{-\theta(x)(x-n\mu)}\left(\mathbb{E}\left( e^{\theta(x) \min\{ X-\mu, x-x^{\varepsilon_0}\} } \right)\right)^n.
\end{align}
Now it follows from Lemma \ref{lem2} that 
\begin{align}\label{e8-1}
	& \mathbb{P}\left(S_n >x, \max_{1\leq i\leq n} X_i \leq x-2x^{\varepsilon_0} \right) \leq e^{-\theta(x)(x-n\mu)} e^{C_*\theta^2(x)n}\nonumber\\
	&=  e^{-\lambda x^{b} +(\log x)^2+\left(C_*\theta^2(x)+ \mu \theta(x) \right)n}.
\end{align}
Since $\ell(x)x^a \gtrsim x^{-|a|-1}$ and  $\log x<n^{1/3}$, the above probability is bounded by
\begin{align}
	& \mathbb{P}\left(S_n >x, \max_{1\leq i\leq n} X_i \leq x-2x^{\varepsilon_0} \right)  \lesssim \mathbb{P}(X>x) e^{ (|a|+1)\log x +(\log x)^2+\left(C_*\theta^2(x)+ \mu \theta(x) \right)n}\nonumber\\
	& \leq \mathbb{P}(X>x) e^{ \left((|a|+1)n^{-2/3} +n^{-1/3}+C_*\theta^2(x)+ \mu \theta(x) \right)n}. 
\end{align}
Therefore, taking $x$ sufficiently large such that 
$(|a|+1)n^{-2/3} +n^{-1/3}+C_*\theta^2(x) + \mu \theta(x) \leq -\log m/2$ for all $n> (\log x)^3$, we conclude that 
\begin{align}\label{e9-1}
		m^{n/2}\mathbb{P}\left(S_n >x, \max_{1\leq i\leq n} X_i \leq x-2x^{\varepsilon_0} \right)&  \lesssim m^{n/2} \mathbb{P}(X>x) e^{-n\log m/2} =  \mathbb{P}(X>x) .
\end{align}
Next, it follows from \eqref{proper-of-tail} and  $\varepsilon_0<1-b$ that
\begin{align}\label{e9-2}
	& m^{n/2} \mathbb{P}\left(S_n >x,  \exists\  1\leq  i\neq j\leq n, X_i, X_j > x-2x^{\varepsilon_0} \right) \nonumber\\
	& \leq m^{n/2} n^2 \left(\mathbb{P}(X> x-2x^{\varepsilon_0})\right)^2 \lesssim   \left(\mathbb{P}(X> x)\right)^2 \leq  \mathbb{P}(X> x).
\end{align}
Combining \eqref{e9-1} and \eqref{e9-2}, we get that when $x$ is large enough, for $(\log x)^3\leq n< x^{(1+b)/2}$,
\begin{align}\label{e9}
	& m^{n/2}\mathbb{P}(T_x^+=n) \leq m^{n/2} \mathbb{P}(S_n>x) \nonumber\\
	&\lesssim   \mathbb{P}(X>x)   + n m^{n/2} \mathbb{P}\left(S_n>x, X_n>x-2x^{\varepsilon_0}, \max_{1\leq j\leq n-1} X_j\leq x-2x^{\varepsilon_0}\right) \nonumber\\
	&\leq    \mathbb{P}(X>x)  + n m^{n/2} \mathbb{P}\left( X_n>x-2x^{\varepsilon_0} \right) \lesssim   \mathbb{P}(X>x)  ,
\end{align}
where the last inequality follows from \eqref{proper-of-tail}.  Combining \eqref{e7}, \eqref{e8} and \eqref{e9}, we complete the proof of (i).

(ii) The case $n=1$ holds trivally, so we consider the case $n\geq 2$. On one hand, from \eqref{e9-2},   for each $n\in \mathbb{N}$,
\begin{align}\label{e11-1}
	\lim_{x\to+\infty} \frac{1}{\mathbb{P}(X>x)}\mathbb{P}\left(T_x^+=n, \exists 1\leq  i\neq j\leq n, X_i, X_j > x-2x^{\varepsilon_0} \right) =0. 
\end{align}
On the other hand, let $x$ be large such that $n< (\log x)^3$, then combining \eqref{proper-of-tail} and \eqref{e6},  it holds that for all $1\leq j\leq n$,
\begin{align}
 	& \lim_{x\to+\infty} \frac{1}{\mathbb{P}(X>x)} \mathbb{P}\left(S_j\in (x-2x^{\varepsilon_0}, x] \right) \nonumber\\
 	&= \lim_{x\to+\infty} \frac{\mathbb{P}\left(S_j> x-2x^{\varepsilon_0} \right)}{\mathbb{P}(X>x-2x^{\varepsilon_0})} \frac{\mathbb{P}\left(X> x-2x^{\varepsilon_0} \right)}{\mathbb{P}(X>x)} -\lim_{x\to+\infty} \frac{\mathbb{P}\left(S_j>x \right)}{\mathbb{P}(X>x)} =0.
\end{align}
Therefore,
\begin{align}
	& \limsup_{x\to+\infty} \frac{1}{\mathbb{P}(X>x)}\mathbb{P}\left(T_x^+=n, \max_{1\leq i\leq n} X_i \leq x-2x^{\varepsilon_0}\right) \nonumber\\
	& \leq  \limsup_{x\to+\infty} \frac{1}{\mathbb{P}(X>x)} \mathbb{P}\left(S_{n-1} \leq x-2x^{\varepsilon_0}, S_n> x,  X_n \leq x-2x^{\varepsilon_0}\right) + \lim_{x\to+\infty} \frac{\mathbb{P}\left(S_{n-1}\in (x-2x^{\varepsilon_0}, x] \right)}{\mathbb{P}(X>x)} \nonumber\\
	& \leq \limsup_{x\to+\infty} \frac{1}{\mathbb{P}(X>x)}\mathbb{P}\left(x-X_n\leq  S_{n-1},  2x^{\varepsilon_0}\leq X_n \leq x-2x^{\varepsilon_0}\right).
\end{align}
Since $S_{n-1}$ and $X_n$ are independent, combining \eqref{e6}, $\ell (x)\lesssim x$ and the fact that $x-X_n \geq 2x^{\varepsilon_0}$ on the event $\{2x^{\varepsilon_0}\leq X_n \leq x-2x^{\varepsilon_0}\}$,  it holds that 
\begin{align}\label{e10}
	& \limsup_{x\to+\infty} \frac{1}{\mathbb{P}(X>x)}\mathbb{P}\left(T_x^+=n, \max_{1\leq i\leq n} X_i \leq x-2x^{\varepsilon_0}\right)\nonumber\\
	& \leq (n-1) \limsup_{x\to+\infty} \frac{1}{\mathbb{P}(X>x)}\mathbb{E}\left(1_{\left\{2x^{\varepsilon_0}\leq X_n \leq x-2x^{\varepsilon_0} \right\}} \mathbb{P}(X \geq x-z)\big|_{z=X_n} \right) \nonumber\\
	&\lesssim  n \limsup_{x\to+\infty} \frac{1}{\mathbb{P}(X>x)} \mathbb{E}\left(1_{\left\{2x^{\varepsilon_0}\leq X_n \leq x-2x^{\varepsilon_0} \right\}}  |x-X_n|^{|a|+1}e^{-\lambda (x-X_n)^b} \right) .
\end{align}
Noticing that $\mathbb{P}(X_n \in (k, k+1])\leq \mathbb{P}(X_n >k)\lesssim k^{|a|+1}e^{-\lambda k^b}$, we have
\begin{align}
	&\mathbb{E}\left(1_{\left\{2x^{\varepsilon_0}\leq X_n \leq x-2x^{\varepsilon_0} \right\}}  |x-X_n|^{|a|+1}e^{-\lambda (x-X_n)^b} \right) \nonumber\\
	&\leq x^{|a|+1} \sum_{k=\lfloor 2x^{\varepsilon_0}\rfloor}^{\lceil x-2x^{\varepsilon_0}\rceil} e^{-\lambda (x-k-1)^b}\mathbb{P}(X_n\in (k, k+1]) \nonumber\\
	&\lesssim x^{|a|+1} \sum_{k=\lfloor 2x^{\varepsilon_0}\rfloor}^{\lceil x-2x^{\varepsilon_0}\rceil}   e^{-\lambda (x-k-1)^b}k^{|a|+1}e^{-\lambda k^{b}}\nonumber\\
	&\leq x^{2|a|+2} \sum_{k=\lfloor 2x^{\varepsilon_0}\rfloor}^{\lceil x-2x^{\varepsilon_0}\rceil}   e^{-\lambda (x-k-1)^b-\lambda k^{b}}.
\end{align}
Since $k^b + (x-k-1)^b\geq  \min_{y\in \{ \lfloor 2x^{\varepsilon_0}\rfloor,  \lceil x-2x^{\varepsilon_0} \rceil \}}(y^b + (x-y-1)^b) \geq (2x^{\varepsilon_0})^b +x^b-1$ for large $x$, we conclude that 
\begin{align}\label{e10-1}
&\mathbb{E}\left(1_{\left\{2x^{\varepsilon_0}\leq X_n \leq x-2x^{\varepsilon_0} \right\}}  |x-X_n|^{|a|+1}e^{-\lambda (x-X_n)^b} \right) \nonumber\\
&\lesssim x^{2|a|+3} e^{-\lambda x^b-\lambda (2x^{\varepsilon_0})^{b}} \lesssim x^{3|a|+4}e^{-\lambda (2x^{\varepsilon_0})^{b}} \mathbb{P}(X>x).
\end{align}
Plugging this back to \eqref{e10} implies that 
\begin{align}\label{e11-2}
	& \lim_{x\to+\infty} \frac{1}{\mathbb{P}(X>x)} \mathbb{P}\left(T_x^+=n, \max_{1\leq i\leq n} X_i \leq x-2x^{\varepsilon_0}\right)= 0. 
\end{align}
Combining \eqref{e11-1} and \eqref{e11-2}, we obtain that
\begin{align}\label{e11-4}
	& \lim_{x\to+\infty} \frac{\mathbb{P}\left(T_x^+=n\right)}{\mathbb{P}(X>x)} \nonumber\\
	&= \lim_{x\to+\infty} \frac{1}{\mathbb{P}(X>x)}\mathbb{P}\left(T_x^+=n, \exists !\ 1\leq j\leq n\ \mbox{such that } \ X_j > x-2x^{\varepsilon_0}\right).
\end{align}
Similar to \eqref{e11-3}, by \eqref{proper-of-tail}, we have the following estimate:
\begin{align}
	 & \lim_{x\to +\infty} \frac{1}{\mathbb{P}(X>x)}\mathbb{P}\left(\exists\  1\leq j\leq n\ \mbox{such that } X_j>x- 2x^{\varepsilon_0}, \max_{\ell\neq j} |X_\ell| > x^{\varepsilon_0} \right) \nonumber\\
	&\leq \lim_{x\to +\infty} \frac{n}{\mathbb{P}(X>x)}\mathbb{P}\left(X_1>x- 2x^{\varepsilon_0}, \max_{\ell\neq 1} |X_\ell| > x^{\varepsilon_0} \right) \nonumber\\
	&=  \lim_{x\to +\infty} \frac{n}{\mathbb{P}(X>x)} \mathbb{P}\left(X>x- 2x^{\varepsilon_0}\right)\mathbb{P}\left(  |X_\ell| > x^{\varepsilon_0} \right)^{n-1}= n \lim_{x\to +\infty} \mathbb{P}\left(  |X_\ell| > x^{\varepsilon_0 }\right)^{n-1}  =0.
\end{align}
Therefore, the limit on the left hand side of \eqref{e11-4} is equal to 
\begin{align}\label{e11-5}
	& \lim_{x\to+\infty} \frac{\mathbb{P}\left(T_x^+=n\right)}{\mathbb{P}(X>x)} \nonumber\\
	& = \lim_{x\to+\infty} \frac{1}{\mathbb{P}(X>x)}\mathbb{P}\left(T_x^+=n, \exists \ 1\leq j\leq n \mbox{ such that } X_j > x-2x^{\varepsilon_0}, \max_{ \ell\neq j} |X_\ell| \leq x^{\varepsilon_0}\right).
\end{align}
If the unique large jump $X_j$ satisfies $j\leq n-1$, then on $\{T_x^+=n \}\cap \{\max_{\ell\neq j} |X_\ell| \leq x^{\varepsilon_0}\}$, it holds that $X_j =S_j- \sum_{\ell=1}^{j-1}X_\ell \leq x+ nx^{\varepsilon_0}$, which together with \eqref{proper-of-tail} implies that 
\begin{align}\label{e11-6}
	&\lim_{x\to+\infty} \frac{1}{\mathbb{P}(X>x)}\mathbb{P}\left(T_x^+=n, \exists\  1\leq j\leq n-1 \mbox{ such that } X_j > x-2x^{\varepsilon_0}, \max_{ \ell\neq j} |X_\ell| \leq x^{\varepsilon_0}\right) \nonumber\\
	&\leq \sum_{j=1}^{n-1}\lim_{x\to+\infty} \frac{1}{\mathbb{P}(X>x)}\mathbb{P}\left( x+ nx^{\varepsilon_0} \geq X_j > x-2x^{\varepsilon_0}\right) =0.
\end{align}
Therefore, combining \eqref{e11-5} and \eqref{e11-6}, we have the following upper bound
\begin{align}\label{eq12-1}
	& \lim_{x\to+\infty} \frac{\mathbb{P}\left(T_x^+=n\right)}{\mathbb{P}(X>x)}  = \lim_{x\to+\infty} \frac{1}{\mathbb{P}(X>x)}\mathbb{P}\left(T_x^+=n,  X_n > x-2x^{\varepsilon_0}, \max_{1\leq  \ell \leq n-1} |X_\ell | \leq  x^{\varepsilon_0}\right) \nonumber\\
	&\leq \lim_{x\to+\infty} \frac{\mathbb{P}\left(X_n > x-2x^{\varepsilon_0}\right)}{\mathbb{P}(X>x)}=1,
\end{align}
where the last equality follows from \eqref{proper-of-tail}. For the lower bound, noticing that on the event $\{ X_n > x+nx^{\varepsilon_0},\max_{ 1\leq \ell \leq n-1} |X_\ell| \leq x^{\varepsilon_0}\}$, we have $\max_{1\leq j\leq n-1} S_j \leq nx^{\varepsilon_0} \leq x$ and $S_n =X_n+S_{n-1}>x$ for large $x$. Therefore,
by \eqref{proper-of-tail},   
\begin{align}\label{eq12-2}
	& \lim_{x\to+\infty} \frac{\mathbb{P}\left(T_x^+=n\right)}{\mathbb{P}(X>x)} \geq \lim_{x\to+\infty} \frac{1}{\mathbb{P}(X>x)}\mathbb{P}\left(X_n > x+nx^{\varepsilon_0},\max_{ 1\leq \ell \leq n-1} |X_\ell| \leq x^{\varepsilon_0} \right) \nonumber\\
	&= \lim_{x\to+\infty} \frac{1}{\mathbb{P}(X>x)}\mathbb{P}\left(X> x+nx^{\varepsilon_0}\right) \mathbb{P}\left( |X| \leq x^{\varepsilon_0} \right)^{n-1}=1.
\end{align}
Combining \eqref{eq12-1} and \eqref{eq12-2} completes the proof of (ii).

\hfill$\Box$



\section{Random walks without sufficient exponential moment}\label{S3}

 In this section, we assume that $X$ satisfies {\bf(H2)}. According to the representation theorem (for example, see \cite[Theorem 1.3.1, p.12]{BGT1987}), there exists $K_0>x_*$  such that 
 \begin{align}\label{Def-of-K-0}
 	0< \inf_{x\in [K_1, K_2]} \ell(x) \leq \sup_{x\in [K_1, K_2]} \ell(x)<\infty \quad\mbox{for all}\quad K_2>K_1\geq K_0.
 \end{align}
 Let $\{X_i, i\in\mathbb{N}\}$ be iid random variables equal in law to $X$. Recall the definitions of $S_n$ and $T_x^+$ in \eqref{Def-of-stopping-time}.  For each $n$, define 
 \begin{align}\label{Change-of-measure}
 	\frac{\mathrm{d} \widetilde{\mathbb{P}}}{\mathrm{d} \mathbb{P}}\bigg|_{\sigma(X_1,...,X_n)}:= \frac{1}{\big(\mathbb{E}(e^{\gamma X})\big)^n} e^{\gamma S_n}.
 \end{align}
 According to {\bf(H2)}, for any $x>x_*$,
 \begin{align}\label{density-X-tildeP}
 	\widetilde{\mathbb{P}}(X>x)= \int_x^\infty \frac{\ell(y)}{\mathbb{E}(e^{\gamma X})} y^a e^{-\lambda y^b}\mathrm{d}y.
 \end{align}
We check here that $(X, \widetilde{\mathbb P})$ satisfies {\bf(H1)}. 

\begin{lemma}\label{lem1}
	There exists some slowly varying function $\widetilde{\ell}$ at $\infty$ such that for all $x>K_0$, 
	\[
		\widetilde{\mathbb{P}}(X>x)=	\widetilde{\ell}(x) x^{a+1-b} e^{-\lambda x^b}.
	\]
	Moreover, $\widetilde{\mathbb{E}}(|X|^k)<\infty$ for all $k>0$ when $b\in (0,1)$ and $\lim_{x\to+\infty} \frac{\widetilde{\ell}(x)}{\ell(x)}= \frac{-(a+1)}{\mathbb{E}(e^{\gamma X})}1_{\{b=0\}}+ \frac{1}{\lambda b \mathbb{E}(e^{\gamma X})}1_{\{b\in (0,1)\}}$.
\end{lemma}
\textbf{Proof: } We divide the proof into two steps.

\noindent
{\bf (Step 1).} In this step, we prove the lemma for $b=0$. From \cite[Proposition 1.5.10, p. 27]{BGT1987}, we have
\begin{align} 
	\lim_{x\to+\infty} \frac{\widetilde{\mathbb{P}}(X>x)}{x^{a+1}\ell(x)}= \frac{-(a+1)}{\mathbb{E}(e^{\gamma X})},
\end{align}
which completes the proof of the lemma by taking $\widetilde{\ell}(x) =x^{-(a+1)}\widetilde{\mathbb{P}}(X>x)$.

\noindent
{\bf (Step 2).} In this step, we prove the lemma for $b\in (0,1)$. Noticing that 
\begin{align}
	\int_x^\infty y^a e^{-\lambda y^b}\mathrm{d} y =  \frac{1}{b}\int_{x^b}^\infty z^{\frac{a+1-b}{b}} e^{-\lambda z}\mathrm{d}z = \frac{x^{a+1-b}e^{-\lambda x^b}}{b}\int_{0}^\infty (1+zx^{-b})^{\frac{a+1-b}{b}} e^{-\lambda z}\mathrm{d}z .
\end{align}
According to dominated convergence theorem,  we get
\begin{align}
	\lim_{x\to+\infty} \frac{\int_x^\infty y^a e^{-\lambda y^b}\mathrm{d} y }{x^{a+1-b}e^{-\lambda x^b}}= \frac{1}{b}\int_0^\infty e^{-\lambda z}\mathrm{d}z = \frac{1}{\lambda b}.
\end{align}
Consequently, with $x$ replaced by $2x$ in the above limit, we obtain 
\begin{align}\label{Limit2}
			\lim_{x\to+\infty} \frac{\int_x^{2x} y^a e^{-\lambda y^b}\mathrm{d} y }{x^{a+1-b}e^{-\lambda x^b}} = \frac{1}{\lambda b}.
\end{align}
According to Lemma \ref{lem-useful-fact} (ii) (with $A=2, \delta=1$), 
\begin{align}\label{e14}
		\widetilde{\mathbb{P}}(X>2x) & =\frac{\ell(x)}{\mathbb{E}(e^{\gamma X})}  \int_{2x}^\infty \frac{\ell(y)}{\ell(x)}y^a e^{-\lambda y^b}\mathrm{d}y\lesssim \ell(x)  \int_{2x}^\infty \frac{y^{a+1}}{x} e^{-\lambda y^b}\mathrm{d}y.
\end{align}
Therefore,  combining  \eqref{e14} and L'H\^opital's rule, we conclude that 
 \begin{align}
 	\limsup_{x\to+\infty} \frac{\widetilde{\mathbb{P}}(X>2x) }{\ell(x) x^{a+1-b} e^{-\lambda x^b}} & \lesssim \lim_{x\to+\infty}\frac{\int_{2x}^\infty y^{a+1} e^{-\lambda y^b}\mathrm{d}y}{x^{a+2-b} e^{-\lambda x^b}} \nonumber\\
 	& = \lim_{x\to+\infty}\frac{-2 (2x)^{a+1} e^{-\lambda (2x)^b} }{ (a+2-b)x^{a+1-b} e^{-\lambda x^b} - \lambda bx^{a+2}e^{-\lambda x^b}} =0.
 \end{align}
 Since $\ell(y)/\ell(x)\to 1$ uniformly for $y\in [x,2x]$ as $x\to+\infty$, combining \eqref{Limit2} and the above limit, it holds that 
 \begin{align}
 		 & \lim_{x\to+\infty} \frac{\widetilde{\mathbb{P}}(X>x) }{\ell(x) x^{a+1-b} e^{-\lambda x^b}}  = 	\lim_{x\to+\infty} \frac{\widetilde{\mathbb{P}}(X>x) - \widetilde{\mathbb{P}}(X>2x) }{\ell(x) x^{a+1-b} e^{-\lambda x^b}}\nonumber\\
 		&= \frac{1}{ \lambda b \mathbb{E}(e^{\gamma X})}\lim_{x\to \infty} \frac{ \int_x^{2x} \ell(y) y^a e^{-\lambda y^b}\mathrm{d}y}{\ell(x) \int_x^{2x} y^a e^{-\lambda y^b}\mathrm{d} y}  =  \frac{1}{ \lambda b \mathbb{E}(e^{\gamma X})}.
 \end{align}
 Now set $\widetilde{\mathbb P}(X>x)= \widetilde{\ell}(x) x^{a+1-b} e^{-\lambda x^b}$, we know that $\widetilde{\ell}$ is a slowly varying function and this completes the proof of the case $b\in (0,1)$.
 
  The check for $\widetilde{\mathbb{E}}(|X|^k)<\infty$ is obvious since the right tail of $(X, \widetilde{\mathbb{P}})$ is stretched exponential and by \eqref{Change-of-measure}, for any $x>0$, 
  \begin{align}\label{left-tail-X}
  	\widetilde{\mathbb{P}}(X<-x) = \frac{1}{\mathbb{E}(e^{\gamma X})}\mathbb{E}(e^{\gamma X}1_{\{X<-x\}}) \leq  e^{-\gamma x}  \frac{1}{\mathbb{E}(e^{\gamma X})} \lesssim e^{-\gamma x}.
  \end{align}
 
 
 
 \hfill$\Box$

 
 
 
 


	\begin{prop}\label{prop2}
	Let $\beta \in (0,1)$ be fixed.
	\begin{itemize}
		\item[(i)] For any $x>K_0$ and $n\in\mathbb{N}$,
		\[
	      \beta^{n/2} \widetilde{\mathbb{E}}\left(e^{-\gamma (S_n-x)}1_{\{T_x^+=n\}}\right)  \lesssim_\beta   \ell(x)x^a e^{-\lambda x^b}.
		\]
		\item[(ii)] For each $n\in \mathbb{N}$,
		\[
		\lim_{x\to+\infty} \frac{ e^{\lambda x^b}}{ \ell(x)x^a } \widetilde{\mathbb{E}}\left(e^{-\gamma (S_n-x)}1_{\{T_x^+=n\}}\right)= \frac{1}{\gamma \mathbb{E}(e^{\gamma X})} \left(1+ \sum_{j=1}^{n-1}\widetilde{\mathbb{E}}\left( \left(1-e^{-\gamma \min_{1\leq i\leq j} S_i}\right)_+\right) \right),
		\]
	\end{itemize}
	here $x_+:= \max\{0,x\}$ and we adapt the convention $\sum_{j=1}^0 = 0$.
\end{prop}


\noindent
\textbf{Proof of Proposition \ref{prop2}  for $b=0$:}  
Since $a<-2$ by {\bf(H2)}, there exists a $\varepsilon\in (1/2,1)$ very close to $1$ such that $(-a)>\frac{2\varepsilon}{2\varepsilon -1}$, which is also equivalent to $(2\varepsilon-1)a+2\varepsilon<0$.   Also, the proof for $n=1$ follows by standard calculation, so we consider the case $n\geq 2$. 

\noindent
{\bf(Step 1)}. In this step, we define a good event $G_n$ and prove a uniform convergence result (see \eqref{e15} below). 
Noticing that $e^{-\gamma (S_n-x)}\leq e^{-\gamma x^{\varepsilon}}$ on the event $\{S_n>x+x^\varepsilon\}$, we have 
\begin{align}\label{e15-1}
	\sup_{n\in \mathbb{N}} \frac{  \beta^{n/2}  }{\ell(x) x^a} \widetilde{\mathbb{E}}\left(e^{-\gamma (S_n-x)}1_{\{T_x^+=n, S_n> x+ x^{\varepsilon}\}}\right) & \leq 	\frac{  e^{-\gamma x^{\varepsilon}} }{\ell(x) x^a} \stackrel{x\to+\infty}{\longrightarrow}0.  
\end{align}
Now define 
\begin{align}\label{def-of-F}
	F_n:= \left\{ \mbox{there exists at most one }i \leq n \mbox{ such that } X_i>x^\varepsilon \right\}.
\end{align}
Then by Lemma \ref{lem1},
\begin{align}\label{e15-2}
	& \frac{  \beta^{n/2}  }{\ell(x) x^a} \widetilde{\mathbb{E}}\left(e^{-\gamma (S_n-x)}1_{\{T_x^+=n, F_n^c\}}\right)  \leq 	\frac{  \beta^{n/2}  }{\ell(x) x^a}  \widetilde{\mathbb{P}}\left(  F_n^c \right)  \leq  \frac{\beta^{n/2} n^2 }{\ell(x) x^a}\widetilde{\mathbb{P}}\left( X>x^\varepsilon \right)^2 \nonumber\\
	& \lesssim_\beta \frac{\ell^2(x^{\varepsilon}) x^{2\varepsilon(a+1)}}{\ell(x) x^a}  = \frac{\ell^2(x^{\varepsilon}) }{\ell(x) } x^{(2\varepsilon-1)a + 2\varepsilon}.
\end{align}
Since $\frac{\ell^2(x^{\varepsilon}) }{\ell(x) } $ is still a slowly varying function at $\infty$ and $(2\varepsilon-1)a + 2\varepsilon<0$, the right hand side of the above display converges to $0$ as $x\to+\infty$ uniformly for all $n\in \mathbb{N}$. 
It follows from \eqref{left-tail-X} that
\begin{align}\label{e15-4}
	&\widetilde{\mathbb{E}}\left(e^{-\gamma (S_n-x)}1_{\{T_x^+=n, \min_{1\leq i\leq n} X_i <-x^{\varepsilon}\}}\right)  \leq  \widetilde{\mathbb{P}}\left( \min_{1\leq i\leq n} X_i< -x^{\varepsilon}\right)\nonumber\\
	& \leq n \widetilde{\mathbb{P}}\left(X<-x^{\varepsilon}\right) \lesssim  ne^{-\gamma x^{\varepsilon}} .
\end{align}
Therefore, we have 
\begin{align}\label{e15-3}
	& \sup_{n\in \mathbb{N}}\frac{ \beta^{n/2} }{\ell(x) x^a}\widetilde{\mathbb{E}}\left(e^{-\gamma (S_n-x)}1_{\{T_x^+=n, \min_{1\leq i\leq n} X_i < -x^{\varepsilon}\}}\right)  \nonumber\\
	&\lesssim  \sup_{n\in \mathbb{N}} n \beta^{n/2}  \frac{e^{-\gamma x^{\varepsilon}}}{\ell(x)x^a} \lesssim_\beta \frac{e^{-\gamma x^{\varepsilon}}}{\ell(x)x^a}\stackrel{x\to+\infty}{\longrightarrow}0. 
\end{align}
In conclusion,  define
\begin{align}
	G_n:= F_n\cap \{S_n\leq x+x^\varepsilon\} \cap \Big\{\min_{1\leq i\leq n} X_i \geq -x^{\varepsilon} \Big\},
\end{align} 
then combining \eqref{e15-1}, \eqref{e15-2} and \eqref{e15-3}, we have 
\begin{align}\label{e15}
	\lim_{x\to+\infty} \sup_{n\in \mathbb{N}}\frac{ \beta^{n/2} }{\ell(x) x^a} \widetilde{\mathbb{E}}\left(e^{-\gamma (S_n-x)}1_{\{T_x^+=n, G_n^c\}}\right) =0. 
\end{align}

\noindent
{\bf(Step 2)}. In this step, we  restrict the expectarion in $G_n$ and prove the proposition. 
 For simplicity we only consider large $x$ and  $n\geq 2$ since   $\inf_{x\in (K_1, K)}\ell(x)x^a >0$ for all $K>K_1$ by \eqref{Def-of-K-0} and the case $n=1$ follows by direct calculation. 
 
 \noindent
 {\bf Case 1: $n  \geq (\log x)^2$.}  We have (noticing that this part is only used in the proof of (i)),
\begin{align}\label{e17}
	 \beta^{n/2} \widetilde{\mathbb{E}}\left(e^{-\gamma (S_n-x)}1_{\{T_x^+=n, G_n\}}\right) \leq  \beta^{n/2}\leq  \beta^{(\log x)^2/2}\lesssim_\beta   \ell(x)x^a.
\end{align}

\noindent
{\bf Case 2: $n  <(\log x)^2$.}  Since on $\{T_x^+=n, \max_{1\leq i\leq n} X_i \leq x^{\varepsilon}\}$, we have  $S_n \leq n x^{\varepsilon}\leq (\log x)^2 x^{\varepsilon} <x$ for large $x$. Therefore, $\{T_x^+=n, \max_{1\leq i\leq n} X_i \leq x^{\varepsilon}\}\cap G_n=\emptyset$, which implies that for large $x$, 
\begin{align}\label{e13}
		  & \widetilde{\mathbb{E}}\left(e^{-\gamma (S_n-x)}1_{\{T_x^+=n, G_n\}}\right)  \nonumber\\
		  &= 	 \widetilde{\mathbb{E}}\left(e^{-\gamma (S_n-x)}1_{\{T_x^+=n, S_n\leq x+x^\varepsilon,  \exists   j\leq n, s.t. X_j>x^\varepsilon,\ \max_{i\neq j} |X_i| \leq x^{\varepsilon}  \}}\right) \nonumber\\
		  & = \sum_{j=1}^n  \widetilde{\mathbb{E}}\left(e^{-\gamma (S_n-x)}1_{\{T_x^+=n, S_n\leq x+x^\varepsilon,   X_j>x^\varepsilon,\ \max_{i\neq j} |X_i| \leq x^{\varepsilon}  \}}\right) .
\end{align}
For each $j\leq n$, $\{S_n>x\}\cap \{\max_{i\neq j} |X_i| \leq x^{\varepsilon}   \}\subset \{X_j \geq x- nx^\varepsilon\}\subset \{X_j >x^\varepsilon\}$. Thus,
\begin{align}\label{e16}
	  & \widetilde{\mathbb{E}}\left(e^{-\gamma (S_n-x)}1_{\{T_x^+=n, G_n\}}\right)   = \sum_{j=1}^n  \widetilde{\mathbb{E}}\left(e^{-\gamma (S_n-x)}1_{\{T_x^+=n, S_n\leq x+x^\varepsilon, \ \max_{i\neq j} |X_i| \leq x^{\varepsilon}  \}}\right) .
\end{align}
To prove (i),   we have 
\begin{align}
	  & \widetilde{\mathbb{E}}\left(e^{-\gamma (S_n-x)}1_{\{T_x^+=n, G_n\}}\right)   \leq \sum_{j=1}^n  \widetilde{\mathbb{E}}\left(e^{-\gamma (S_n-x)}1_{\{ x<S_n\leq x+x^\varepsilon, \ \max_{i\neq j} |X_i| \leq x^{\varepsilon}  \}}\right) \nonumber\\
	  & = n \widetilde{\mathbb{E}}\left(1_{\{ \max_{1\leq i\leq n-1} |X_i| \leq x^{\varepsilon}  \}} \widetilde{\mathbb{E}} \left(e^{-\gamma (X-y)}1_{\{ y< X\leq y+x^\varepsilon\}}\right)\Big|_{y=x-S_{n-1}}\right). 
\end{align}
Since $|S_{n-1}|\leq nx^\varepsilon \leq (\log x)^2 x^{\varepsilon} $ on $\{ \max_{1\leq i\leq n-1} |X_i| \leq x^{\varepsilon}  \}$, we have
\begin{align}\label{e18}
	& \frac{ \beta^{n/2}}{\ell(x) x^a} \widetilde{\mathbb{E}}\left(e^{-\gamma (S_n-x)}1_{\{T_x^+=n, G_n\}}\right) \leq 	 \frac{ n\beta^{n/2} }{\ell(x) x^a} \sup_{|y-x|\leq (\log x)^2 x^{\varepsilon}} \widetilde{\mathbb{E}} \left(e^{-\gamma (X-y)}1_{\{ y< X\leq y+x^\varepsilon\}}\right)\nonumber\\
	& \lesssim_\beta  \frac{  1}{\ell(x) x^a}\sup_{|y-x|\leq (\log x)^2 x^{\varepsilon}}  \int_0^{x^\varepsilon} \ell(z+y) (z+y)^a e^{-\gamma z}\mathrm{d}z \nonumber\\
	& \lesssim 1. 
\end{align}
Combining \eqref{e15}, \eqref{e17} and \eqref{e18}, we complete the proof of (i).

Now we are ready to prove (ii). Since $\max_{1\leq k\leq j-1}S_k \leq n x^\varepsilon <x$ for large $x$ on $\{\max_{i\neq j} |X_i| \leq x^{\varepsilon}\}$, we have for each $n\in \mathbb{N}$ and $1\leq j\leq n-1$, 
\begin{align}\label{e19}
	& \widetilde{\mathbb{E}}\left(e^{-\gamma (S_n-x)}1_{\{T_x^+=n, S_n\leq x+x^\varepsilon, \ \max_{i\neq j} |X_i| \leq x^{\varepsilon}  \}}\right) \nonumber\\
	& =  \widetilde{\mathbb{E}}\left(e^{-\gamma (S_n-x)}1_{\{ X_j \leq x-\max_{j+1\leq k\leq n}(S_{k-1}-X_j),\  X_j + (S_n-X_j) \in (x,x+x^\varepsilon], \ \max_{i\neq j} |X_i| \leq x^{\varepsilon}  \}}\right). 
\end{align}
If $j=n$, we  remove the first event $\{ X_j \leq x-\max_{j+1\leq k\leq n}(S_{k-1}-X_j)\}$. 
 Set $Y=x-( S_n-X_j)$, $U=S_n- \max_{j\leq k\leq n-1} S_k$ for $1\leq j \leq n-1$ and $U=\infty$ for $j=n$. According to the independence between $X_j$ and $(U, Y)$, we have for all $1\leq j\leq n$, 
 \begin{align}
 		& \widetilde{\mathbb{E}}\left(e^{-\gamma (S_n-x)}1_{\{T_x^+=n, S_n\leq x+x^\varepsilon, \ \max_{i\neq j} |X_i| \leq x^{\varepsilon}  \}}\right) \nonumber\\
 	& =  \widetilde{\mathbb{E}}\left(e^{-\gamma (X_j-Y)} 1_{\{ X_j \leq U+Y,\  X_j  \in (Y,Y+x^\varepsilon], \ \max_{i\neq j} |X_i| \leq x^{\varepsilon} , U>0 \}}\right)\nonumber\\
 	& = \widetilde{\mathbb{E}}\left(1_{\{\max_{i\neq j} |X_i| \leq x^{\varepsilon} \}} 1_{\{U>0\}}\widetilde{\mathbb{E}}\left(e^{-\gamma (X_j-y)}1_{\{ X_j \leq u+y,\  X_j  \in (y,y+x^\varepsilon] \}}\right)\Big|_{y=Y, u=U}\right).
 \end{align}
We also mention here that on the event $\{\max_{i\neq j} |X_i| \leq x^{\varepsilon} \}$,  $ |Y-x| \leq n x^\varepsilon \leq (\log x)^2 x^\varepsilon$.  Since  the law of $U= S_n- \max_{j\leq k\leq n-1} S_k$ is independent of $x$ and that  uniformly for all $|y-x| \leq (\log x)^2 x^\varepsilon$,
\begin{align}
	 & \widetilde{\mathbb{E}}\left(e^{-\gamma (X_j-y)}1_{\{ X_j \leq u+y,\  X_j  \in (y,y+x^\varepsilon] \}}\right)  = \frac{1}{\mathbb{E}(e^{\gamma X})}\int_0^{\min\{x^{\varepsilon}, u\} } \ell(z+y) (z+y)^a e^{-\lambda z} \mathrm{d}z,
\end{align}
 taking $x\to+\infty$ in the above limit implies that uniformly on $|y-x| \leq (\log x)^2 x^\varepsilon$,
\begin{align}
	& \frac{ 1 }{\ell(x)x^a}\widetilde{\mathbb{E}}\left(e^{-\gamma (X_j-y)}1_{\{ X_j \leq u+y,\  X_j  \in (y,y+x^\varepsilon] \}}\right)= \frac{1}{\mathbb{E}(e^{\gamma X})}\int_0^{\min\{ x^{\varepsilon}, u\} }\frac{ \ell(z+y)}{\ell (x)} \left(\frac{z+y}{x}\right)^a e^{-\lambda z} \mathrm{d}z \nonumber\\
	 & \stackrel{x\to+\infty}{\longrightarrow} \frac{1}{\mathbb{E}(e^{\gamma X})} \int_0^ue^{-\lambda z}\mathrm{d}z.
\end{align}
Noticing that $S_n- \max_{j\leq k\leq n-1} S_k$ is equal in law to $\min_{1\leq i\leq n-j} S_i$, we conclude that 
\begin{align}
	& \lim_{x\to+\infty} \frac{1}{\ell(x)x^a}\widetilde{\mathbb{E}}\left(e^{-\gamma (S_n-x)}1_{\{T_x^+=n, G_n\}}\right)   \nonumber\\
	&= \frac{1}{\mathbb{E}(e^{\gamma X})} \left(\frac{1}{\gamma}+ \sum_{j=1}^{n-1}\widetilde{\mathbb{E}}\left(1_{\{\min_{1\leq i\leq j} S_i>0 \}}\int_0^{\min_{1\leq i\leq j} S_i} e^{-\gamma z}\mathrm{d} z\right) \right).
\end{align}
Therefore, (ii) follows from \eqref{e15} and the above display. We are done. 

\hfill$\Box$


\noindent

\textbf{Proof of Proposition \ref{prop2}  for $b\in (0,1)$:}  The proof is similar to the case $b=0$ with more delicate details. Since the proof for $n=1$ follows by standard calculation,   we consider the case $n\geq 2$.  Recall that $\varepsilon_0= \frac{1}{2}\min\{ b, 1-b\}$. 


\begin{thebibliography}{99}
\bibitem[BGT1987]{BGT1987} Bingham N. H., Goldie C. M. and Teugels J. L.: \emph{Regular variation, Encyclopedia Math. Appl., 27}. Cambridge University Press, Cambridge, 1989
\bibitem[DDS2008]{DDS2008} Denisov, D., Dieker, A. B. and Shneer., V.: Large deviations for random walks under subexponentiality: The big-jump domain. \emph{Ann. Probab.} \textbf{36} (5) (2008) 1946--1991.
\end{thebibliography}
\end{document}
